Si assumano tutti gli operazionali ideali ed operanti in regime lineare. Si considerino i valori
delle resistenze pari a quelli nominali.

\paragraph{Richiesta (a)}
Mostrare sulla base di semplici argomenti che i diodi D1 e D2 non possono essere
simultaneamente nello stesso stato di polarizzazione (se cioè uno è in conduzione l’altro deve
essere interdetto e viceversa).
\paragraph{Risposta (a)}
\begin{quote}
  ``Entrambi i diodi sono entrambi in conduzione?''
\end{quote}
Assumendo senza perdita di generalit\'a una tensione d'ingresso \textit{positiva}
sul canale $V_S$ allora la tensione fra le resistenze e i diodi dovranno essere di segno
\textit{concorde} siccome gli elementi resistivi non possiedono reattanza. Assumendo per
assurdo che entrambi i diodi siano in conduzione si ottiene uno schema di amplificatore
\textit{invertente}, dunque la tensione all'uscita dell'amplificatore $V_{out}$ avr\'a
segno \textit{discorde} rispetto tensione d'ingresso $V_S$. Il diodo ``al di sotto''
dell'ingresso non-invertente non pu\'o essere messo in conduzione perch\'e all'ingresso
avr\'a potenziale \textit{negativo} e all'uscita potenziale \textit{positivo}, una
contraddizione (un argomento analogo viene fatto assumendo la tensione d'ingresso \textit{negativa}).

\begin{quote}
  ``Entrambi i diodi sono entrambi interdetti?''
\end{quote}
Se si assumesse invece che nessuno dei diodi \'e in conduzione allora la tensione all'uscita
dell'amplificatore \'e nulla, quindi nuovamente almeno uno fra i diodi entra in conduzione
(perch\'e la tensione d'ingresso \'e non nulla), un'altra contraddizione.
\textit{Almeno uno} dei diodi dovr\'a essere in conduzione.

\paragraph{Richiesta (b)}
Alla luce del precedente argomento si giustifichino le identità tra le semi-onde di VA, VB, VS
di cui al punto 1b.
\paragraph{Risposta (b)}
Siccome un diodo alla volta pu\'o entrare in conduzione le possibili tensioni all'uscite
di $V_A$ e $V_B$ sono rispettivamente
\begin{equation*}
  \begin{gathered}
    V_S > 0 \implies \frac{V_S - V_{-U1}}{R_1} = \frac{V_{-U1} - V_{+U3}}{R_2} \\
    V_S < 0 \implies \frac{V_S - V_{-U1}}{R_1} = \frac{V_{-U1} - V_{+U4}}{R_3}
  \end{gathered}
\end{equation*}
siccome $R_1 = R_2 = R_3$ e $V_{-U1} = V_{+U1} = 0$ allora
\begin{equation*}
  \begin{gathered}
    V_S > 0 \implies
    \begin{cases}
      V_{+U3} = -V_S \\
      V_{+U4} = 0
    \end{cases} \\
    V_S < 0 \implies
    \begin{cases}
      V_{+U3} = 0 \\
      V_{+U4} = -V_S
    \end{cases}
  \end{gathered} \implies
  V_{+U3} + V_{+U4} = -V_S
\end{equation*}
inoltre $V_{+U3}$ e $V_{+U4}$ sono entrambi delle tensioni d'ingresso a degli \textit{inseguitori}
di tensione con rispettive uscite $V_A$ e $V_B$, infine dovr\'a valere (scritto in maniera elegante
ma forse meno leggibile)
\begin{equation*}
  V_A = -\max\{V_S, 0\} \qquad V_B = -\min\{V_S, 0\} \qquad V_A + V_B = -V_S
\end{equation*}

\paragraph{Richiesta (c)}
Determinare la funzione di trasferimento del secondo stadio del circuito (con ingressi VA e
VB ed uscita VOUT, verificando che si tratti di un amplificatore differenziale a guadagno
complesso).

\paragraph{Richiesta (d)}
Calcolare il modulo del guadagno di modo differenziale atteso nel limite di continua ($f\to 0$)
e ad 1kHz.

\paragraph{Richiesta (e)}
Calcolare il valore atteso per il rapporto tra il livello di continua di VOUT e l’ampiezza di VS
e lo si confronti con la pendenza della retta di best-fit ottenuta al punto 1d.

\paragraph{Richiesta (f)}
Si dia una stima del valore atteso percentuale dell’ampiezza della prima
armonica di VOUT, ovvero del rapporto tra la sua escursione ed il suo valor medio, e lo si
metta in relazione con le misura di ripple di cui al punto 1b.
